
All experimental data analyzed in this study was previously collected and described in several publications (\cite{duszkiewicz,Siegle2021-hs, Siegel2015-cz, iblreproducibility}).

\subsection{Shape Metrics}

\subsubsection{Procrustes distance between neural populations of unequal sizes}

In the main text, we introduced the Procrustes distance calculation as an orthogonal alignment between two neural population recordings, represented by $N \times M$ matrices $\mbX_i$ and $\mbX_j$ (see eq.~\ref{eq:procrustes-distance}).
This assumes that each neural recording consists of $N$ neurons, because the transformation matrix $\mbQ$ must be square (i.e. $N \times N$) in order to be orthogonal.

This assumption is unrealistic because most experimental recordings will yield different numbers of neurons.
However, the Procrustes distance can be easily generalized to this scenario while preserving the essential geometric intuition of our main narrative.
We demonstrate this in a few ways below.

For the rest of this subsection we consider $\mbX_i$ be an $N_i \times M$ matrix and $\mbX_j$ be an $N_j \times M$ matrix where $N_i \neq N_j$.
We assume that the rows of $\mbX_i$ and $\mbX_j$ have been preprocessed to have mean zero.

\textbf{Option 1 --- PCA preprocessing.}
The first strategy we consider is using PCA to project each dataset onto an $N$-dimensional subspace.
If the projection onto this lower-dimensional space captures a large fraction of variance in each dataset, then we intuitively expect the Procrustes distance within this low-dimensional space to be a reasonable measure of representational similarity.

Formally, let $\mbP_i$ be an $N \times N_i$ orthonormal projection matrix onto the top $N$ principal component subspace of $\mbX_i$ (i.e. top $N$ eigenvectors of $\mbX_i \mbX_i^\top$).
Similarly, let $\mbP_j$ be a $N \times N_j$ orthonormal projection matrix onto the top $N$ principal component subspace of $\mbX_j$.
Then define $\widetilde{\mbX}_i = \mbP_i \mbX_i / \Vert \mbP_i \mbX_i \Vert_F$ and $\widetilde{\mbX}_j = \mbP_j \mbX_j / \Vert \mbP_j \mbX_j \Vert_F$ as normalized projections of the data onto the top $N$ principal components.
We suggest to normalize by $\Vert \mbP_i \mbX_i \Vert_F$ and $\Vert \mbP_j \mbX_j \Vert_F$ to account for the fact that there will likely be a mismatch in overall variance between the two datasets.
That is, the ``size'' of the manifolds in the projected space may be quite different, simply because one system had more neurons and thus more variance to begin with.
Another reasonable normalization strategy would be to set $\widetilde{\mbX}_i = \mbP_i \mbX_i / \sqrt{N_i}$ and $\widetilde{\mbX}_j = \mbP_j \mbX_j / \sqrt{N_j}$.

If the projection onto a lower, $N$-dimensional space captures a large fraction of variance in each dataset, then we intuitively expect the Procrustes distance between $\mbP_i \mbX_i$ and $\mbP_j \mbX_j$ to be a reasonable measure of representational similarity (since we are only neglecting low-variance dimensions).
Thus, to retain maximal variance in each dataset we can set $N = \min(N_i, N_j)$ when comparing a pair of neural systems.
When making comparisons among $K$ systems, we must take $N = \min(N_1, \dots, N_K)$ to preserve symmetry and the triangle inequality.

\textbf{Option 2 --- Generalize the definition of Procrustes distance.}
While the approach outlined above is intuitive, it is essentially a preprocessing trick that circumvents deeper issues.
Furthermore, it is hard to justify in cases where non-trivial amounts of variance are lost by projecting the data down into $N$ dimensions.
We now outline a more elegant approach to think about the problem and show that it largely coincides with the intuitive PCA-driven approach described above.

We begin by stating a classical mathematical result that the Procrustes distance between two matrices can be written in a closed form~\parencite{schonemann1966generalized}.
Specifically, the squared distance can be written as:
\begin{equation}
\label{eq:procrustes-closed-form-1}
\min_{\mbQ \in \cO(N)} \Vert \mbQ \mbX_i - \mbX_j \Vert_F^2 = \Tr[ \mbX_i \mbX_i^\top ] + \Tr[ \mbX_j \mbX_j^\top ] - 2 \Vert \mbX_i \mbX_j^\top \Vert_*
\end{equation}
where $\Vert \mbX_i \mbX_j^\top \Vert_*$ denotes the sum of the singular values of $\mbX_i \mbX_j^\top$ and is called the \textit{nuclear norm} or \textit{Schatten 1-norm}.
We provide a derivation of this classic result in \hl{appendix}.

The equation above describes the case where $N_i = N_j$ (i.e. the two systems have equal numbers of neurons), but the expression on the right hand side can be calculated even when $N_i \neq N_j$ because the nuclear norm is well-defined even when $\mbX_i \mbX_j^\top$ is rectangular.


\subsection{Downstream Analyses}

Shape metrics provide a strategy for computing a $K \times K$ distance matrix, $\mbD$, that captures pairwise dissimilarity across a collection of $K$ neural systems.
Once this distance matrix is obtained, there are a number of analyses that can be immediately performed using off-the-shelf tools.

\subsubsection{Multidimensional Scaling (MDS)}

We use a public implementation of metric MDS from \texttt{scikit-learn} \parencite{scikit}.
Briefly, given a $K \times K$ matrix $\mbD$ of pairwise distances, MDS uses iterative parameter updates in an attempt to solve the following nonconvex optimization problem
\begin{equation}
\underset{\mbu_1, \dots, \mbu_K}{\text{minimize}} \quad \sum_{i=1}^K \sum_{j=i+1}^K \left ( \mbD_{ij} -  \Vert \mbu_i - \mbu_j \Vert_2 \right )^2
\end{equation}
where $\mbu_1, \dots, \mbu_K \in \reals^d$ are vectors within a $d$-dimensional embedding space.
Once obtained, these vectors can be used for any downstream analysis (e.g. running PCA to visualize each neural system in a low-dimensional space; see Figures \ref{fig:synth}b-c, \ref{fig:hdcells}c, \ref{fig:allen}c). 
However, MDS comes at the cost of introducing distortion into the metric space. 
We quantify the distortion of distance $\mbD_{ij}$ as
\begin{equation}
\max \left ( \frac{\mbD_{ij}}{\Vert \mbu_i - \mbu_j \Vert_2}, \frac{\Vert \mbu_i - \mbu_j \Vert_2}{\mbD_{ij}} \right ) 
\end{equation}
and report the median distortion across all distance values.

\subsection{Datasets and Analyses}

\subsubsection{Simulated Dataset (Fig. \ref{fig:synth})}

We randomly generated $K = 40$ sets of tuning curves (corresponding to different ``networks'' or ``subjects'').
Each simulated network consisted of $N = 100$ neurons with tuning curves given by a 2-parameter family of functions:
\begin{equation}
\label{eq:methods-tuning-curve-shape}
f(\theta) = \exp \left ( -\kappa \cos (\theta - \mu)\right )
\end{equation}
Intuitively, $\mu \in [-\pi, \pi)$ describes the peak of the tuning curve and $\kappa \geq 0$ inversely controls the width of the tuning.
Under this model, the peak firing rate of each neuron is equal to one.
Note that \cref{eq:methods-tuning-curve-shape} is proportional to the density of a von Mises distribution, for which $\kappa$ is called the concentration parameter and $\mu$ is called the location parameter.

To generate the dataset, we specify a value of $\kappa$ and a distribution over $\mu$ for each subject and sample $N=100$ random tuning curves.
The concentration parameter for each subject, $\kappa$, was sampled uniformly between 1 and 4; all tuning curves from each subject were assigned the same $\kappa$ parameter.
The location parameter for each tuning curve, $\mu$, was sampled from a von-mises distribution with mean zero and a concentration parameter $\gamma$ that was drawn uniformly between 0 and 3 for each subject.
Intuitively, a simulated network will roughly have uniformly spaced tuning curves when the concentration parameter controlling the distribution of $\mu$ is close to 0.
Conversely, a simulated network will have tuning curves that are ``bunched'' around $\theta=0$ when this concentration parameter is close to 3.

We computed the acuity of each subject by taking the logarithm of the expected Fisher Information (FI) across neurons evaluated at $\theta=0$.
Specifically, under the assumption of isotropic Gaussian noise with variance $\sigma^2$, the expected FI is given by
\begin{equation}
\label{eq:fish-info}
I = \expect \left [ \frac{f'(0)^2}{\sigma^2} \right ] 
\end{equation}
where the expected value is taken over randomly sampled tuning curves, $f$.
Note that in our simulated data we only get to observe $N=100$ tuning curves to compute the Procrustes distances between network pairs.
This simulates a realistic scenario where we only get to experimentally record a small fraction of the neural circuit underlying behavioral performance.
Nonetheless, if $N$ is large enough and the shape of the neural representation is smooth/low-dimensional, we can expect the empirical Procrustes distance to do a reasonable job of capturing the ``true'' geometry of the system.

Applying some basic algebraic manipulations to \cref{eq:fish-info}, we obtain
\begin{equation}
I = \frac{\kappa^2}{\sigma^2} \int \sin^2(\mu) \exp(-2\kappa \cos(\mu)) p(\mu) \textrm{d}\mu
\end{equation}
where $p(\mu)$ is the density of a von-mises distribution with mean $\theta=0$ and concentration parameter $\gamma$.
Although some further simplifications are possible, the form of the resulting integral does not appear to have a simple analytic solution.
Thus, we numerically computed it to high precision to obtain $\log I$, the simulated ``acuity'' score for each subject.
We set $\sigma^2 = 1$ in this calculation; this choice is of little consequence since it only introduces an overall scaling into the simulation.


% \subsection{Dissimilarity matrix}

% All analyses performed here started by constructing a dissimilarity matrix. All entries in this matrix consisted of a distance in shape space, but depending on the analyzes, this could be the distance between subjects, conditions, or regions. Unless stated otherwise, dissimilarity matrices where cross-validated (see below).

% \subsection{Permutation metric}

% The group of permutation matrices $\mathcal{P}(N)$ is a subgroup of $\cO(N)$ (see also eq. \ref{eq:generalized-euclidean-shape-metric}), therefore a distance calculated on neurons matched one-to-one across networks is also a rigorous metric. To evaluate this metric, we optimize to the set of neuron permutations that maximally align two given networks. As shown in \cite{Williams2020}, this can be reformulated as a problem in combinatorial optimization known as the linear assignment problem. 

% \subsection{Supervised and unsupervised learning in shape space}
% We leveraged the rigorous metric space defined by shape analyses to perform different downstream analyzes, notably regression, clustering, and PCA. For these analyzes, we used off-the-shelf libraries available in scikit-learn (\cite{scikit}). For the ordinal regression performed in (Fig. \ref{fig:miller}d), we used threshold-based ordinal logistic regression (LogisticAT from pythonhosted.org/mord). For hierarchical clustering performed in Fig. \ref{allen}, we used scipy \parencite{scipy}. When regressing shape spaces against behavioral statistics (e.g. Fig. \ref{fig:hdcells}c, right or Fig. \ref{fig:ibl}d) we always did so using different partitions as targets and features (see Cross-valdiation, below)  Before performing dimensionality reduction of a dissimilarity matrix using PCA, we embedded its vectors in shape space into euclidean space using multi-dimensional scaling, an approach that has been shown to minimize distortions (\cite{Williams2021}).

%  \subsection{Cross-validation}
% We employed different types of cross-validation before trial-averaging. Typically, we averaged independent trials/time-points, leading to independent partitions across trials/time-points. Additionally, we also performed a more conservative approach in which we partitioned also across neurons of the same animal, leading to completely independent dataset partitions from the same animal (Fig. \ref{fig:hdcells}). This process was repeated n=1000 times, unless stated otherwise.

% \begin{center}
% \begin{tikzpicture}

% % Matrix outline
% \draw[thick] (0, 0) rectangle (8, 6);

% % Horizontal partition (across Trials)
% \draw[thick, dashed] (0, 3) -- (8, 3);

% % Vertical partition (across Neurons)
% \draw[thick, dashed] (4, 0) -- (4, 6);

% % Labels
% \node at (4, 6.3) {Trials};
% \node[rotate=90] at (-0.3, 3) {Neurons};

% % Labels for partitions
% \node at (2, 4.5) {training partition};
% \node at (2, 1.5) {ignored};

% \node at (6, 4.5) {ignored};
% \node at (6, 1.5) {testing partition};
% \end{tikzpicture}
% \end{center}


\subsubsection{Head direction coding in mice (Fig. \ref{fig:hdcells})}

We analyzed two datasets made public by \textcite{duszkiewicz}.
Spiking activity from cells in the mouse postsubiculum (PoSub) was collected using silicone probes during open-field exploration. 
The first dataset consisted of $31$ sessions/subjects and a total of $2691$ neurons. The second dataset is similar to the first one, but included a salient landmark that was rotated 90$^{\circ}$  back and forth every 200s for a total of 16 rotations. 

In both datasets, spikes were counted in bins of 1$^{\circ}$ of the head direction and then bins were smoothed with 12$^{\circ}$  s.d. Gaussian filters, as done in the original publication \parencite{duszkiewicz}.
Each partition for cross-validation was then denoised by keeping only the first 20 principal components (see Fig. \ref{fig:hdcells}a).  

To perform leave-one-subject out decoding, we aligned all subjects to a common reference. The common reference was chosen randomly and did not change the results. We aligned each subject manifold using n=20 randomly chosen, but equally spaced, head direction bins. Importantly, we did not include these angles or $\pm$6 bins around them in the training set for head-direction decoding.

\subsubsection{Allen brain observatory dataset (Fig. \ref{fig:allen})}

We analyzed the ``natural movie three'' stimulus set of the Neuropixels---Visual Coding dataset in the Allen Brain Observatory (\cite{Siegle2021-hs}).
Data were obtained and analyzed through the AllenSDK python package.
We pooled data from 32 recording sessions, which is the total number of sessions employing the \textit{brain observatory 1.1} stimulus set.
Spike counts were binned at the video presentation frame rate ($30^{-1}$ second time bins), averaged across ten trials, and smoothed by convolving with a Gaussian kernel (standard deviation of 20 time bins).
This resulted in $M=3600$ time bins of activity measurements for each neuron.
Neurons were then pooled across sessions and separated into subregions using the most fine-grained anatomical regions defined by the Allen Brain Atlas.
We excluded brain regions with fewer than $50$ neurons and additionally excluded subregions with the acronyms \texttt{\{MB, alv, ccs, dhc, fp, or\}} which represent either white matter or midbrain neurons that were not assignable to any specific nucleus.
This resulted in a total of $K=55$ brain regions.
To denoise these datasets and speed up calculations, we projected each dataset onto the top $N=50$ principal components before analysis.

\subsubsection{Nonhuman primate visually guided decision-making task (Fig. \ref{fig:miller})}

Acute recordings from 7 regions of the monkey cortex were performed using single tungsten/parylene electrodes (\cite{Siegel2015-cz}). Monkeys were trained to perform a context-dependent decision-making task. The stimulus consisted of dots moving in a random direction and of a random color. 

We computed condition-averaged responses (n=64 conditions, 4 colors x 4 directions x 4 contextual cues; \cite{Siegel2015-cz}) from $N=7542$  multiunits recorded from 4 brain regions from subject R, and 7 regions from subject P (Fig. \ref{fig:miller}A). For each condition, we averaged two independent sets of trials separately (folds, in Fig. \ref{fig:miller}). We repeated this process 10 times, leading to 10 datasets consisting of 2 independent folds of 11 brain regions. Finally, to avoid across-area differences to be dominated by different dataset sizes, we subsampled each dataset to n=80 neurons and reduced their dimensionality to N=5 using PCA, leading to a data tensor of dimensions 2 x 11 x 5 for each condition that we analyse using shape metrics. All the analyses were repeated 100 times and we report the averages. 

\subsubsection{International Brain laboratory dataset (Fig. \ref{fig:ibl})}
\hl{add autism dataset}
We analysed a subset of the IBL dataset published in \cite{iblreproducibility}. Briefly, the IBL labs conducted Neuropixels recordings from the same brain regions in head-fixed mice during a standardized decision-making task, where the mice reported the perceived position of a visual grating. The experiment followed a uniform pipeline across all labs, including standardized methods for surgery, behavioral training, recordings, histology, and data processing. We focused here exclusively on behavior and spike counts recordings. 

In addition to the original quality control, we only analyzed subjects from which at least $n = 100$ neurons were recorded and with at least $n = 40$ trials per condition were recorded. \hl{Moreover, we only focused on data collected during balanced blocks (i.e. in which the probability of rewarded left and right responses was the same). In total, we analysed data from $n = 56$ subjects.}

